\section{Discussion}
\label{sec:discussion}

\subsection{Case Study: \texttt{count\_up\_down-1}}

The canonical VGuide win is \texttt{count\_up\_down-1}, the motivating example from
\Cref{sec:intro}. The program has a single loop in which \texttt{x} is incremented and
\texttt{y} is decremented by one on each iteration, starting from \texttt{x = 0} and
\texttt{y = n}. After the loop, the property \texttt{assert(y == 0)} must hold. The
proof requires the relational invariant $x + y = n$, which is an affine constraint
relating two loop variables to a parameter.

Under stock CEGAR, the refinement sequence is as follows. The first spurious CE
witnesses the loop executing zero times; interpolation adds $\{y = n\}$ to $\Pi$.
The second CE witnesses one iteration; interpolation adds $\{y = n - 1 \wedge x = 1\}$.
Each subsequent CE witnesses one more iteration and adds one more conjunction. After
\CUDStockRefs{} such refinements, the 300\,s budget expires at \CUDStockWall{}\,s wall
time and CPAChecker returns \textsc{Unknown}.

VGuide fires on the first spurious CE. The context pack includes the loop source,
the block formula for the one-iteration CE, and the CFA loop-head node label. From
this context, \LLMName{} proposes nine candidate predicates, including \texttt{x + y == n},
\texttt{y == n - x}, \texttt{x >= 0}, \texttt{y <= n}, and \texttt{x <= n}. The
three-tier cascade accepts five of the nine (two are redundant by L3; two fail L2 due
to a variable name variant). The five accepted predicates are injected at the loop head.
On the next CEGAR iteration, the abstract transformer with the enriched $\Pi$ can
represent the exact invariant $x + y = n$. CEGAR converges in \CUDVGuideRefs{}
further refinements and returns \textsc{True} at \CUDVGuideWall{}\,s wall time---a
reduction of 95\% in wall time and 97\% in refinement count.

\subsection{Failure Mode: \texttt{nested-3}}

The clearest regression against stock CEGAR is \texttt{nested-3}, a program with a
two-level nested loop. The outer loop is a \texttt{while} loop that resets a state
variable \texttt{st} to 1 at the loop head. The inner \texttt{for} loop conditionally
sets \texttt{st} to 0 on a subset of iterations. The outer loop's safety invariant
requires tracking whether the inner loop's conditional branch has been taken, making
it sensitive to the exact Boolean structure of \texttt{st}.

VGuide fires on the first spurious CE and the LLM proposes predicates including
\texttt{st == 0} and \texttt{st == 1} for the outer loop head. Both predicates are
syntactically valid, in scope, and satisfiable (L3 checks each predicate individually
for satisfiability, not mutual exclusivity). They are therefore accepted and injected.
However, the joint injection of $\{st = 0, st = 1\}$ at the outer loop head causes
the predicate abstraction to split the loop's abstract state into four combinations,
two of which ($st = 0 \wedge st = 1$) are inconsistent but not immediately pruned.
This expands the CEGAR search space by a factor proportional to the predicate set size
and changes the refinement trajectory: CEGAR now explores abstract states that would
not have been generated from interpolation alone.

The result is that \texttt{nested-3} requires 42 refinements and times out at
\textsc{Unknown}, compared to stock CEGAR which produces \textsc{True} in 2 refinements
in 1.3\,s. This failure mode reveals the key risk of predicate injection: adding
semantically correct but contextually noisy predicates can increase rather than
decrease the CEGAR search space. An L3 extension that checks mutual consistency across
injected predicates (rather than each in isolation) could mitigate this failure mode
at the cost of additional SMT queries.

\subsection{Limitations}

\paragraph{Online LLM dependency.} VGuide calls an external LLM API during
verification. In the standard SV-COMP competition setting, participating tools execute
in a network-isolated environment; no external API access is permitted. VGuide in its
current form cannot be submitted as a standalone SV-COMP competition tool. Mitigations
include pre-analysis LLM calls (querying the LLM before the competition isolation
window), locally hosted models, or a pre-computed predicate cache keyed by program
fingerprint.

\paragraph{FALSE-path blindness.} VGuide improves only the \textsc{True} (safety
proof) count. The \textsc{False} (bug-finding) count is identical across all VGuide
and stock runs. Predicate hints accelerate proof convergence by providing better
abstract-domain coverage, but they do not provide counterexample witnesses for genuine
bugs. Extending VGuide with a BUG-mode predicate strategy---injecting predicates that
help CEGAR refine toward the reachable error region---is a natural direction for
improving \textsc{False} performance.

\paragraph{Sound predicates, noisy hints.} The three-tier validation cascade guarantees
that accepted predicates are semantically valid, but it does not guarantee that they
are \emph{useful} for the proof at hand. Noisy predicates---those that are correct but
irrelevant to the current property---consume refinement budget by increasing the
abstract state space without moving CEGAR closer to termination. The \texttt{nested-3}
regression is an instance of this. Future work on predicate relevance filtering, using
heuristics derived from the property formula or the CE structure, may reduce the
fraction of noisy accepted predicates.
